\documentclass{article}

% Adds the url command
\usepackage{hyperref}
% To remove paragraph indentation
\usepackage{parskip}
% Adds line cut to urls
\usepackage{xurl}
\usepackage{multirow}
\usepackage{algorithm}
\usepackage{algpseudocode}

% Update these:
\title{Lab 11 Report}
\date{Due: 04/28/2026 @ 11:00 AM}
\author{Ryan Wilson}
\begin{document}
\maketitle
\newpage

\section{Complexity}

Using the Gauss-Jordan elimination method, the solver requires $\mathcal{O}(n^3)$ operations for n unknowns. This means that the overall complexity for this project is $\mathcal{O}(n^3)$. Parsing the input file is $\mathcal{O}(n)$, and constructing the matrix is $\mathcal{O}(n^2)$, so they are negligible when compared to the solver.

\section{Loop Analysis}

The loops with the most runtime happen in the Gauss-Jordan elimination process, which uses three nested loops over the number of unknowns. The matrix construction uses nested loops over nodes and edges, giving a complexity of $\mathcal{O}(n^2)$. One loop invariant is that, during elimination, columns become zero except at the diagonal. Another loop invariant is that, during back substitution, solved variables stay unchanged. Non-dominant loops include the parsing and file output loops.

\section{Sources}
AERE361\_LabManual\_11.pdf\\
\url{https://www.geeksforgeeks.org/compiler-design/introduction-of-parsing-ambiguity-and-parsers-set-1/}\\
\url{https://www.geeksforgeeks.org/linux-unix/bash-scripting-introduction-to-bash-and-bash-scripting/}\\


\end{document}
